\IfFileExists{stacks-project.cls}{%
\documentclass{stacks-project}
}{%
\documentclass{amsart}
}

% For dealing with references we use the comment environment
\usepackage{verbatim}
\newenvironment{reference}{\comment}{\endcomment}
%\newenvironment{reference}{}{}
\newenvironment{slogan}{\comment}{\endcomment}
\newenvironment{history}{\comment}{\endcomment}

% For commutative diagrams we use Xy-pic
\usepackage[all]{xy}

% We use 2cell for 2-commutative diagrams.
\xyoption{2cell}
\UseAllTwocells

% We use multicol for the list of chapters between chapters
\usepackage{multicol}

% This is generally recommended for better output
\usepackage{lmodern}
\usepackage[T1]{fontenc}

% For cross-file-references
\usepackage{xr-hyper}

% Package for hypertext links:
\usepackage{hyperref}

% For any local file, say "hello.tex" you want to link to please
% use \externaldocument[hello-]{hello}
\externaldocument[introduction-]{introduction}
\externaldocument[conventions-]{conventions}
\externaldocument[sets-]{sets}
\externaldocument[categories-]{categories}
\externaldocument[topology-]{topology}
\externaldocument[sheaves-]{sheaves}
\externaldocument[sites-]{sites}
\externaldocument[stacks-]{stacks}
\externaldocument[fields-]{fields}
\externaldocument[algebra-]{algebra}
\externaldocument[brauer-]{brauer}
\externaldocument[homology-]{homology}
\externaldocument[derived-]{derived}
\externaldocument[simplicial-]{simplicial}
\externaldocument[more-algebra-]{more-algebra}
\externaldocument[smoothing-]{smoothing}
\externaldocument[modules-]{modules}
\externaldocument[sites-modules-]{sites-modules}
\externaldocument[injectives-]{injectives}
\externaldocument[cohomology-]{cohomology}
\externaldocument[sites-cohomology-]{sites-cohomology}
\externaldocument[dga-]{dga}
\externaldocument[dpa-]{dpa}
\externaldocument[sdga-]{sdga}
\externaldocument[hypercovering-]{hypercovering}
\externaldocument[schemes-]{schemes}
\externaldocument[constructions-]{constructions}
\externaldocument[properties-]{properties}
\externaldocument[morphisms-]{morphisms}
\externaldocument[coherent-]{coherent}
\externaldocument[divisors-]{divisors}
\externaldocument[limits-]{limits}
\externaldocument[varieties-]{varieties}
\externaldocument[topologies-]{topologies}
\externaldocument[descent-]{descent}
\externaldocument[perfect-]{perfect}
\externaldocument[more-morphisms-]{more-morphisms}
\externaldocument[flat-]{flat}
\externaldocument[groupoids-]{groupoids}
\externaldocument[more-groupoids-]{more-groupoids}
\externaldocument[etale-]{etale}
\externaldocument[chow-]{chow}
\externaldocument[intersection-]{intersection}
\externaldocument[pic-]{pic}
\externaldocument[weil-]{weil}
\externaldocument[adequate-]{adequate}
\externaldocument[dualizing-]{dualizing}
\externaldocument[duality-]{duality}
\externaldocument[discriminant-]{discriminant}
\externaldocument[derham-]{derham}
\externaldocument[local-cohomology-]{local-cohomology}
\externaldocument[algebraization-]{algebraization}
\externaldocument[curves-]{curves}
\externaldocument[resolve-]{resolve}
\externaldocument[models-]{models}
\externaldocument[functors-]{functors}
\externaldocument[equiv-]{equiv}
\externaldocument[pione-]{pione}
\externaldocument[etale-cohomology-]{etale-cohomology}
\externaldocument[proetale-]{proetale}
\externaldocument[relative-cycles-]{relative-cycles}
\externaldocument[more-etale-]{more-etale}
\externaldocument[trace-]{trace}
\externaldocument[crystalline-]{crystalline}
\externaldocument[spaces-]{spaces}
\externaldocument[spaces-properties-]{spaces-properties}
\externaldocument[spaces-morphisms-]{spaces-morphisms}
\externaldocument[decent-spaces-]{decent-spaces}
\externaldocument[spaces-cohomology-]{spaces-cohomology}
\externaldocument[spaces-limits-]{spaces-limits}
\externaldocument[spaces-divisors-]{spaces-divisors}
\externaldocument[spaces-over-fields-]{spaces-over-fields}
\externaldocument[spaces-topologies-]{spaces-topologies}
\externaldocument[spaces-descent-]{spaces-descent}
\externaldocument[spaces-perfect-]{spaces-perfect}
\externaldocument[spaces-more-morphisms-]{spaces-more-morphisms}
\externaldocument[spaces-flat-]{spaces-flat}
\externaldocument[spaces-groupoids-]{spaces-groupoids}
\externaldocument[spaces-more-groupoids-]{spaces-more-groupoids}
\externaldocument[bootstrap-]{bootstrap}
\externaldocument[spaces-pushouts-]{spaces-pushouts}
\externaldocument[spaces-chow-]{spaces-chow}
\externaldocument[groupoids-quotients-]{groupoids-quotients}
\externaldocument[spaces-more-cohomology-]{spaces-more-cohomology}
\externaldocument[spaces-simplicial-]{spaces-simplicial}
\externaldocument[spaces-duality-]{spaces-duality}
\externaldocument[formal-spaces-]{formal-spaces}
\externaldocument[restricted-]{restricted}
\externaldocument[spaces-resolve-]{spaces-resolve}
\externaldocument[formal-defos-]{formal-defos}
\externaldocument[defos-]{defos}
\externaldocument[cotangent-]{cotangent}
\externaldocument[examples-defos-]{examples-defos}
\externaldocument[algebraic-]{algebraic}
\externaldocument[examples-stacks-]{examples-stacks}
\externaldocument[stacks-sheaves-]{stacks-sheaves}
\externaldocument[criteria-]{criteria}
\externaldocument[artin-]{artin}
\externaldocument[quot-]{quot}
\externaldocument[stacks-properties-]{stacks-properties}
\externaldocument[stacks-morphisms-]{stacks-morphisms}
\externaldocument[stacks-limits-]{stacks-limits}
\externaldocument[stacks-cohomology-]{stacks-cohomology}
\externaldocument[stacks-perfect-]{stacks-perfect}
\externaldocument[stacks-introduction-]{stacks-introduction}
\externaldocument[stacks-more-morphisms-]{stacks-more-morphisms}
\externaldocument[stacks-geometry-]{stacks-geometry}
\externaldocument[moduli-]{moduli}
\externaldocument[moduli-curves-]{moduli-curves}
\externaldocument[examples-]{examples}
\externaldocument[exercises-]{exercises}
\externaldocument[guide-]{guide}
\externaldocument[desirables-]{desirables}
\externaldocument[coding-]{coding}
\externaldocument[obsolete-]{obsolete}
\externaldocument[fdl-]{fdl}
\externaldocument[index-]{index}

% Theorem environments.
%
\theoremstyle{plain}
\newtheorem{theorem}[subsection]{Theorem}
\newtheorem{proposition}[subsection]{Proposition}
\newtheorem{lemma}[subsection]{Lemma}

\theoremstyle{definition}
\newtheorem{definition}[subsection]{Definition}
\newtheorem{example}[subsection]{Example}
\newtheorem{exercise}[subsection]{Exercise}
\newtheorem{situation}[subsection]{Situation}

\theoremstyle{remark}
\newtheorem{remark}[subsection]{Remark}
\newtheorem{remarks}[subsection]{Remarks}

\numberwithin{equation}{subsection}

% Macros
%
\def\lim{\mathop{\mathrm{lim}}\nolimits}
\def\colim{\mathop{\mathrm{colim}}\nolimits}
\def\Spec{\mathop{\mathrm{Spec}}}
\def\Hom{\mathop{\mathrm{Hom}}\nolimits}
\def\Ext{\mathop{\mathrm{Ext}}\nolimits}
\def\SheafHom{\mathop{\mathcal{H}\!\mathit{om}}\nolimits}
\def\SheafExt{\mathop{\mathcal{E}\!\mathit{xt}}\nolimits}
\def\Sch{\mathit{Sch}}
\def\Mor{\mathop{\mathrm{Mor}}\nolimits}
\def\Ob{\mathop{\mathrm{Ob}}\nolimits}
\def\Sh{\mathop{\mathit{Sh}}\nolimits}
\def\NL{\mathop{N\!L}\nolimits}
\def\CH{\mathop{\mathrm{CH}}\nolimits}
\def\proetale{{pro\text{-}\acute{e}tale}}
\def\etale{{\acute{e}tale}}
\def\QCoh{\mathit{QCoh}}
\def\Ker{\mathop{\mathrm{Ker}}}
\def\Im{\mathop{\mathrm{Im}}}
\def\Coker{\mathop{\mathrm{Coker}}}
\def\Coim{\mathop{\mathrm{Coim}}}

% Boxtimes
%
\DeclareMathSymbol{\boxtimes}{\mathbin}{AMSa}{"02}

%
% Macros for moduli stacks/spaces
%
\def\QCohstack{\mathcal{QC}\!\mathit{oh}}
\def\Cohstack{\mathcal{C}\!\mathit{oh}}
\def\Spacesstack{\mathcal{S}\!\mathit{paces}}
\def\Quotfunctor{\mathrm{Quot}}
\def\Hilbfunctor{\mathrm{Hilb}}
\def\Curvesstack{\mathcal{C}\!\mathit{urves}}
\def\Polarizedstack{\mathcal{P}\!\mathit{olarized}}
\def\Complexesstack{\mathcal{C}\!\mathit{omplexes}}
% \Pic is the operator that assigns to X its picard group, usage \Pic(X)
% \Picardstack_{X/B} denotes the Picard stack of X over B
% \Picardfunctor_{X/B} denotes the Picard functor of X over B
\def\Pic{\mathop{\mathrm{Pic}}\nolimits}
\def\Picardstack{\mathcal{P}\!\mathit{ic}}
\def\Picardfunctor{\mathrm{Pic}}
\def\Deformationcategory{\mathcal{D}\!\mathit{ef}}
\usepackage{cleveref}
\crefname{originallemma}{original lemma}{original lemmas}
\Crefname{originallemma}{Original lemma}{Original lemmas}
\makeatletter
\newlabel{lemma-application-nice-K-injective}{{08D4}{source}{}{sourcepointer-1}{}}
\newlabel{lemma-application-nice-K-injective@cref}{{[originallemma][0][]08D4}{[1][0][]source}{}{}{}}
\makeatother
\crefname{originallemma}{original lemma}{original lemmas}
\Crefname{originallemma}{Original lemma}{Original lemmas}
\makeatletter
\newlabel{morphisms-lemma-affine-equivalence-modules}{{01SB}{source}{}{sourcepointer-2}{}}
\newlabel{morphisms-lemma-affine-equivalence-modules@cref}{{[originallemma][0][]01SB}{[1][0][]source}{}{}{}}
\makeatother
\crefname{originallemma}{original lemma}{original lemmas}
\Crefname{originallemma}{Original lemma}{Original lemmas}
\makeatletter
\newlabel{modules-lemma-pullback-quasi-coherent}{{01BG}{source}{}{sourcepointer-3}{}}
\newlabel{modules-lemma-pullback-quasi-coherent@cref}{{[originallemma][0][]01BG}{[1][0][]source}{}{}{}}
\makeatother
\crefname{originallemma}{original lemma}{original lemmas}
\Crefname{originallemma}{Original lemma}{Original lemmas}
\makeatletter
\newlabel{coherent-lemma-relative-affine-vanishing}{{01XC}{source}{}{sourcepointer-4}{}}
\newlabel{coherent-lemma-relative-affine-vanishing@cref}{{[originallemma][0][]01XC}{[1][0][]source}{}{}{}}
\makeatother
\title{For an affine morphism the pushforward is an equivalence of quasi-coherent derived categories}
\author{The Stacks Project Authors}
\date{stacks-project:perfect.tex:local-snapshot-20261012; GNU FDL 1.2 or later}
\begin{document}
\maketitle
\setcounter{subsection}{1}
\noindent\textit{Editorial note.} All delivered mathematical text is the author's own, copied from the pinned Stacks Project chapter file perfect.tex; no statement or proof step was rewritten, reordered or removed. The counted span is the author environment and its own proof environment. References inside the span are handed to the assembler as navigation stubs that carry the canonical tag of the cited tag only; whether each of them is load-bearing is recorded on the stub itself.\par\medskip
\noindent\textit{Allowed standard prerequisites (editorial).}\begin{itemize}
\item The derived pushforward is computed on an object by choosing a K-injective representative and taking the pushforward of that complex. used at proof lines 940-946. Defines the functor whose equivalence with the pushforward to the structure sheaf is proved.
\item A K-injective complex of quasi-coherent modules on an affine scheme is described by the corresponding statement so that the pushforward can be compared sheaf by sheaf. invoked at proof lines 946-960. Gives the compatibility needed to view the pushforward as a functor to modules over the pushforward of the structure sheaf.
\item On affine schemes the equivalence between modules over the structure sheaf and over the pushforward of the structure sheaf is the standard one, and pullbacks of quasi-coherent modules remain quasi-coherent. invoked at proof lines 960-985. Reduces the statement to the affine case and shows the comparison functor is an equivalence there.
\item Higher direct images of quasi-coherent modules vanish for an affine morphism, which makes the comparison with the derived pushforward an isomorphism. invoked at proof lines 985-989. Identifies the composition of the equivalence with the pushforward to the structure sheaf with the derived pushforward.
\end{itemize}\medskip
\ifdefined\hypertarget\hypertarget{sourcepointer-1}{}\fi
\noindent\textit{Source reference.} Original lemma 08D4: sources/stacks-perfect/pinned-original-source.tex, source lines 375--375. Invoked at proof lines 946-958 to compare the pushforward of a K-injective complex with the corresponding complex of modules; the reduction to affines that follows is inside the counted span.\par\medskip
\ifdefined\hypertarget\hypertarget{sourcepointer-2}{}\fi
\noindent\textit{Source reference.} Original lemma 01SB: sources/stacks-chow/tags.txt, source lines 1--1. Cited at proof lines 966-975 for the affine case of the comparison; the compatibility with the derived functors is inside the counted span.\par\medskip
\ifdefined\hypertarget\hypertarget{sourcepointer-3}{}\fi
\noindent\textit{Source reference.} Original lemma 01BG: sources/stacks-chow/tags.txt, source lines 1--1. Used at proof lines 976-982 to keep the functors within the quasi-coherent categories; the conclusion drawn is inside the counted span.\par\medskip
\ifdefined\hypertarget\hypertarget{sourcepointer-4}{}\fi
\noindent\textit{Source reference.} Original lemma 01XC: sources/stacks-chow/tags.txt, source lines 1--1. Invoked at proof lines 985-988 to identify the composite functor with the derived pushforward; that identification is the last step inside the counted span.\par\medskip
\begin{lemma}
\label{lemma-affine-morphism-equivalence}
Let $f : X \to Y$ be an affine morphism of schemes. Then $f_*$ induces
an equivalence
$$
\Phi : D_\QCoh(\mathcal{O}_X) \longrightarrow D_\QCoh(f_*\mathcal{O}_X)
$$
whose composition with $D_\QCoh(f_*\mathcal{O}_X) \to D_\QCoh(\mathcal{O}_Y)$
is $Rf_* : D_\QCoh(\mathcal{O}_X) \to D_\QCoh(\mathcal{O}_Y)$.
\end{lemma}

\begin{proof}
Recall that $Rf_*$ is computed on an object $K \in D_\QCoh(\mathcal{O}_X)$
by choosing a K-injective complex $\mathcal{I}^\bullet$ of
$\mathcal{O}_X$-modules representing $K$ and taking $f_*\mathcal{I}^\bullet$.
Thus we let $\Phi(K)$ be the complex $f_*\mathcal{I}^\bullet$
viewed as a complex of $f_*\mathcal{O}_X$-modules.
Denote $g : (X, \mathcal{O}_X) \to (Y, f_*\mathcal{O}_X)$ the
obvious morphism of ringed spaces. Then $g$ is a flat morphism of
ringed spaces (see below for a description of the stalks) and
$\Phi$ is the restriction of $Rg_*$ to $D_\QCoh(\mathcal{O}_X)$.
We claim that $Lg^*$ is a quasi-inverse. First, observe that
$Lg^*$ sends $D_\QCoh(f_*\mathcal{O}_X)$ into $D_\QCoh(\mathcal{O}_X)$
because $g^*$ transforms quasi-coherent modules into quasi-coherent
modules (Modules, Lemma \ref{modules-lemma-pullback-quasi-coherent}).
To finish the proof it suffices to show that
the adjunction mappings
$$
Lg^*\Phi(K) = Lg^*Rg_*K \to K
\quad\text{and}\quad
M \to Rg_*Lg^*M = \Phi(Lg^*M)
$$
are isomorphisms for $K \in D_\QCoh(\mathcal{O}_X)$ and
$M \in D_\QCoh(f_*\mathcal{O}_X)$. This is a local question, hence
we may assume $Y$ and therefore $X$ are affine.

\medskip\noindent
Assume $Y = \Spec(B)$ and $X = \Spec(A)$. Let
$\mathfrak p = x \in \Spec(A) = X$ be a point mapping to
$\mathfrak q = y \in \Spec(B) = Y$. Then
$(f_*\mathcal{O}_X)_y = A_\mathfrak q$ and $\mathcal{O}_{X, x} = A_\mathfrak p$
hence $g$ is flat. Hence $g^*$ is exact and $H^i(Lg^*M) = g^*H^i(M)$
for any $M$ in $D(f_*\mathcal{O}_X)$.
For $K \in D_\QCoh(\mathcal{O}_X)$ we see that
$$
H^i(\Phi(K)) = H^i(Rf_*K) = f_*H^i(K)
$$
by the vanishing of higher direct images
(Cohomology of Schemes, Lemma \ref{coherent-lemma-relative-affine-vanishing})
and Lemma \ref{lemma-application-nice-K-injective} (small detail omitted).
Thus it suffice to show that
$$
g^*g_*\mathcal{F} \to \mathcal{F}
\quad\text{and}\quad
\mathcal{G} \to g_*g^*\mathcal{G}
$$
are isomorphisms where $\mathcal{F}$ is
a quasi-coherent $\mathcal{O}_X$-module and $\mathcal{G}$ is
a quasi-coherent $f_*\mathcal{O}_X$-module. This follows from
Morphisms, Lemma \ref{morphisms-lemma-affine-equivalence-modules}.
\end{proof}

\end{document}
